Place the measured base endpoints at $(-b/2,0)$ and $(b/2,0)$ and the apex at $(d,h)$. Subtracting the two squared-distance equations gives the project's geometric diagnostic:
\begin{equation}
d=\frac{\ell^2-r^2}{2b},\qquad \phi_{\rm app}=\arctan(d/h).
\label{eq:res-apex}
\end{equation}
Using $b=80.30$~mm, $\ell=56.20$~mm, $r=56.90$~mm and the nominal $h=40$~mm gives $d=-0.493$~mm and $\phi_{\rm app}=-0.706^\circ$. The angle is conditional on the assumed height and a shear interpretation; it is not an independent measurement of rail squareness. Both an ideal square and a rhombus can have equal side lengths, which is why diagonal or coordinate measurements would add information absent from the square-side test.

